\section{Results}\label{sec:results}
本章では、前章で述べた手法に基づき得られた実験結果について詳述する。評価指標として精度（Accuracy）および損失（Loss）を用い、モデルの性能を定量的に分析した。

\subsection{評価指標の分析}
表\ref{tab:metrics}に示す通り、学習データおよびテストデータにおいて高い精度を達成した。特にテストデータにおける精度は0.95であり、モデルの汎化性能が十分に確保されていることが確認できる。

\begin{table}[htbp]
\centering
\begin{tabular}{lrr}
\toprule
Metric & Train & Test \\
\midrule
Accuracy & 0.98 & 0.95 \\
Loss & 0.02 & 0.05 \\
\bottomrule
\end{tabular}
\caption{評価指標の比較}
\label{tab:metrics}
\end{table}

\subsection{大規模データセットにおける挙動}
また、異なる実行環境（Run A, B, C）における性能比較を行った。表\ref{tab:wide}は、各環境下での結果をまとめたものである。環境による性能差は限定的であり、本手法の堅牢性が示唆された。

\begin{table*}[htbp]
\centering
\begin{tabular}{lrrr}
\toprule
Run & A & B & C \\
\midrule
R1 & 1 & 2 & 3 \\
R2 & 4 & 5 & 6 \\
\bottomrule
\end{tabular}
\caption{実行環境別の性能比較}
\label{tab:wide}
\end{table*}

\subsection{考察の要点}
以上の結果から、提案手法は特定の条件下において安定した性能を発揮することが明らかとなった。次章では、これらの結果が持つ意義と、残された課題について議論する。
